% =====================================================================
%  2bat-ud2-7.tex  —  SESSIÓ 7: Consolidació: traduir i representar restriccions
% =====================================================================
\horatitol{Sessió 7 — Consolidació: traduir i representar restriccions}%
{Reforç abans de l'optimització: caçar els errors típics i practicar el camí de l'enunciat a la regió factible.}

\subsection{Idea clau}

Avui no hi ha matèria nova: consolidem les dues habilitats que sostenen tota la
unitat —\textbf{traduir} un enunciat a un sistema d'inequacions i \textbf{representar}
la regió factible— caçant els errors típics. L'objectiu és que el camí «enunciat
$\rightarrow$ variables $\rightarrow$ inequacions $\rightarrow$ regió $\rightarrow$
vèrtexs» surti amb seguretat.

\begin{idea}
\textbf{Els errors que cacem avui:}
\begin{itemize}[leftmargin=1.4em, itemsep=2pt]
  \item \textbf{Tipus de línia} equivocat: contínua per a $\le,\ge$; discontínua per a
        $<,>$.
  \item \textbf{Semiplà equivocat}: cal comprovar un punt de prova (sovint l'origen).
  \item \textbf{Oblidar la no-negativitat}: $x\ge 0$ i $y\ge 0$ sempre hi són en
        contextos reals.
  \item \textbf{Vèrtex mal calculat}: cal resoldre bé el sistema de les dues vores.
\end{itemize}
\textbf{Rutina segura:} defineix què és $x$ i què és $y$, tradueix cada condició,
comprova un punt per fixar el semiplà, i calcula els vèrtexs resolent sistemes.
\end{idea}

\subsection{Exemples}

\begin{exemple}[caça l'error: semiplà equivocat]
Algú representa $x+y\le 8$ i ombreja el costat \emph{de dalt} de la recta (lluny de
l'origen). \textbf{És incorrecte:} amb el punt de prova $(0,0)$, $0+0=0\le 8$ és
\textbf{cert}, així que el semiplà solució és el \emph{de l'origen} (a sota de la
recta). \textbf{Correcció:} ombrejar el costat de l'origen.
\end{exemple}

\begin{exemple}[caça l'error: oblidar la no-negativitat]
Algú representa només $x+2y\le 10$ i dóna com a regió tot el semiplà, incloent-hi
parts amb $x<0$ o $y<0$. \textbf{És incorrecte} en un context real: no es poden fer
quantitats negatives. Cal afegir $x\ge 0$ i $y\ge 0$, que confinen la regió al
\textbf{primer quadrant}. \textbf{Correcció:} la regió és només la part del primer
quadrant.
\end{exemple}

\subsection{Exercicis resolts}

\begin{exresolt}[de l'enunciat a la regió]
Una entitat ofereix $x$ places de matí i $y$ de tarda. Cada plaça de matí ocupa $1$
monitor i cada plaça de tarda, $2$; hi ha com a màxim $10$ monitors. A més, en total
no es poden oferir més de $8$ places. Tradueix, representa i dóna els vèrtexs.
\solucio
\textbf{Variables:} $x$ places de matí, $y$ de tarda. \textbf{Condicions:}
\[
x+2y\le 10 \ (\text{monitors}),\quad x+y\le 8 \ (\text{places}),\quad x\ge 0,\ y\ge 0.
\]
\textbf{Vèrtexs:} els eixos donen $(0,0)$. Sobre $y=0$, mana la més petita de
$x\le 10$ i $x\le 8$: $(8,0)$. Creuament de $x+y=8$ i $x+2y=10$: restant, $y=2$,
$x=6$, vèrtex $(6,2)$. Sobre $x=0$: mana la més petita de $y\le 5$ i $y\le 8$:
$(0,5)$.

\begin{center}
\begin{tikzpicture}
\begin{axis}[udaxis, width=8.4cm, height=5.4cm,
    xlabel={\small matí ($x$)}, ylabel={\small tarda ($y$)},
    xmin=0, xmax=11, ymin=0, ymax=7, xtick={0,2,4,6,8,10}, ytick={0,2,4,6}]
  \fill[udblue!18] (axis cs:0,0) -- (axis cs:8,0) -- (axis cs:6,2) -- (axis cs:0,5) -- cycle;
  \addplot[udblue, domain=0:8, samples=2]{8-x};
  \addplot[udnavy, domain=0:10, samples=2]{(10-x)/2};
  \addplot[udnavy, only marks, mark=*, mark size=1.4pt] coordinates {(0,0)(8,0)(6,2)(0,5)};
\end{axis}
\end{tikzpicture}
\end{center}
\resposta{$x+2y\le 10$, $x+y\le 8$, $x,y\ge 0$; vèrtexs $(0,0)$, $(8,0)$, $(6,2)$, $(0,5)$.}
\end{exresolt}

\begin{exresolt}[corregir un vèrtex mal calculat]
Algú diu que les rectes $x+y=8$ i $x+2y=10$ es creuen a $(4,4)$. Comprova-ho i
corregeix si cal.
\solucio
A $(4,4)$: $x+y=8$ (d'acord) però $x+2y=4+8=12\neq 10$. No és el creuament. Resolem
bé: restant les equacions, $(x+2y)-(x+y)=10-8\Rightarrow y=2$; llavors $x+2=8\Rightarrow
x=6$. El creuament correcte és $(6,2)$.
\resposta{És incorrecte; el creuament és $(6,2)$.}
\end{exresolt}

\subsection{Exercicis per fer}

\noindent\textit{Itinerari: 1–3 són la base (traducció guiada); 4–6 són ampliació.}

\begin{exercicis}
\begin{exlist}
  \item \textbf{(Caça l'error)} Digues què està malament en cada cas i corregeix-ho:
        \begin{enumerate}[label=\alph*), leftmargin=1.6em, topsep=2pt]
          \item Per a $x+y<5$ s'ha dibuixat la recta frontera \textbf{contínua}.
          \item Per a $2x+y\ge 6$ s'ha ombrejat el costat de l'origen.
        \end{enumerate}
  \item Tradueix a un sistema d'inequacions: «una entitat fa $x$ tallers i $y$ xerrades;
        en total com a màxim $12$ activitats, i no més de $8$ xerrades; cap quantitat
        negativa».
  \item Representa el sistema $\;x+y\le 6,\;y\le 4,\;x\ge 0,\;y\ge 0\;$ i dóna'n els
        vèrtexs.
  \item \textbf{(Ampliació)} Una associació reparteix $x$ lots de menjar i $y$ lots
        d'higiene. Cada lot de menjar costa $2$ unitats de pressupost i cada lot
        d'higiene, $1$; el pressupost és $10$. A més, com a màxim $4$ lots de menjar.
        Tradueix, representa i dóna tots els vèrtexs.
  \item \textbf{(Ampliació)} Calcula el vèrtex on es creuen $3x+y=12$ i $x+y=6$.
  \item \textbf{(Ampliació)} Per al sistema $\;x+y\le 7,\;x+3y\le 15,\;x\ge 0,\;y\ge 0$,
        troba els quatre vèrtexs de la regió factible.
\end{exlist}
\end{exercicis}

% ---------------------------------------------------------------------
\begin{idea}
\textbf{Full d'autoavaluació.} Marca amb sinceritat com et trobes, de cara a la prova:
\begin{center}
\renewcommand{\arraystretch}{1.4}
\begin{tabular}{p{8.4cm} c c c}
\toprule
\textbf{Sé fer\dots} & \textbf{Sol} & \textbf{Amb ajuda} & \textbf{Encara no} \\
\midrule
Definir les variables a partir d'un enunciat & & & \\
Traduir cada condició a una inequació & & & \\
Triar el tipus de línia i el semiplà correctes & & & \\
Recordar la no-negativitat ($x,y\ge 0$) & & & \\
Calcular els vèrtexs resolent sistemes & & & \\
\bottomrule
\end{tabular}
\end{center}
\end{idea}

% ---------------------------------------------------------------------
\fullsolucions{Sessió 7}
\begin{solucions}
\begin{sollist}
  \item (a) Amb $<$ (estricte) la línia ha de ser \textbf{discontínua}, no contínua. \quad
        (b) Prova $(0,0)$: $0\ge 6$ és fals, així que el semiplà és el \emph{contrari}
        a l'origen; estava ombrejat al revés.
  \item $x+y\le 12$, \;$y\le 8$, \;$x\ge 0$, \;$y\ge 0$.
  \item Vèrtexs $(0,0)$, $(6,0)$, $(2,4)$ i $(0,4)$. (Creuament de $x+y=6$ i $y=4$:
        $x=2$.)
  \item $2x+y\le 10$ (pressupost), $x\le 4$ (lots de menjar), $x\ge 0$, $y\ge 0$.
        Vèrtexs: $(0,0)$, $(4,0)$, $(4,2)$ i $(0,10)$. (A $x=4$: $y\le 2$;
        creuament amb l'eix $y$ a $(0,10)$.)
  \item Restant: $(3x+y)-(x+y)=12-6\Rightarrow 2x=6\Rightarrow x=3$; $y=6-3=3$. Vèrtex
        $(3,3)$.
  \item Vèrtexs $(0,0)$, $(7,0)$, $(3,4)$ i $(0,5)$. (Creuament de $x+y=7$ i
        $x+3y=15$: $2y=8\Rightarrow y=4$, $x=3$.)
\end{sollist}
\end{solucions}
